\begin{table}[H]\centering
\caption{\textbf{Private cesareans with no recorded clinical indication}}
\label{tab:no_indication}\small
\small\setlength{\tabcolsep}{4pt}\resizebox{\ifdim\width>\linewidth \linewidth\else\width\fi}{!}{%
\begin{tabular}{cc}
\toprule
Year & Share with no indication ICD-10 code (\%) \\
\midrule
2015 & 87.2 \\
2016 & 88.2 \\
2017 & 87.9 \\
2018 & 88.1 \\
2019 & 87.9 \\
2020 & 89.0 \\
2021 & 89.2 \\
2022 & 90.0 \\
2023 & 90.8 \\
2024 & 91.2 \\
\bottomrule
\end{tabular}
}
\\[2pt]\footnotesize\textit{Notes:} TISS private cesarean deliveries, 2015--2024. A delivery is coded ``no indication'' when the primary diagnosis (ICD-10 code) is a delivery-outcome code (O80--O84) or blank, rather than an ICD-10 code recording a recognized cesarean indication (malpresentation, disproportion, placental or fetal complications, obstructed labor, etc.). Diagnosis coding in claims is incomplete, so this describes recorded indications, not clinical necessity.
\end{table}
